% ============================================================
%  Section 4.5 – Tests du χ²
% ============================================================

\subsection{Tests du $\chi^2$}

% ── T4 ──────────────────────────────────────────────────────────────────────
\paragraph{T4 -- Conformité de V12 (moyenne au bac) à une loi uniforme.}

\textbf{Hypothèses.}
$H_0$ : la distribution de V12 sur les $k = 4$ modalités est uniforme
(chaque modalité équiprobable, $p_i = 1/4$) ;
$H_1$ : la distribution n'est pas uniforme.

\textbf{Effectifs théoriques.}
$e_i = n \times p_i = 67/4 = 16{,}75$ pour chaque modalité.

\textbf{Conditions} (§\,3.4.2) :
$n = 67 \geq 50$ {\checkmark}, $\min(e_i) = 16{,}75 \geq 5$ {\checkmark}.

\begin{table}[H]
  \centering
  \caption{Effectifs observés et théoriques -- V12}
  \label{tab:chi2_T4}
  \small
  \begin{tabular}{@{}lrrrr@{}}
    \toprule
    Modalité      & 8--11 & 12--14 & 15--17 & 18--20 \\
    \midrule
    Observé $o_i$ & 13    & 39     & 12     & 3      \\
    Théorique $e_i$ & 16{,}75 & 16{,}75 & 16{,}75 & 16{,}75 \\
    \bottomrule
  \end{tabular}
\end{table}

\textbf{Critère.}
\[
  \chi^2_\text{obs}
  = \sum_{i=1}^{4} \frac{(o_i - e_i)^2}{e_i}
  = \frac{(13-16{,}75)^2}{16{,}75} + \frac{(39-16{,}75)^2}{16{,}75}
  + \frac{(12-16{,}75)^2}{16{,}75} + \frac{(3-16{,}75)^2}{16{,}75}
  \approx \mathbf{43{,}03}
\]

ddl $= k - 1 = 3$ ; seuil $\chi^2_{3\,;\,5\%} = 7{,}81$.

\textbf{Conclusion.}
$\chi^2_\text{obs} = 43{,}03 > 7{,}81$ : on \textbf{rejette $H_0$}.
La distribution des niveaux de bac n'est pas uniforme : les modalités « 12--14 »
représentent à elles seules 58,2\,\% de l'échantillon.
La population étudiée est donc homogène en termes de niveau scolaire initial.

% ── T5 ──────────────────────────────────────────────────────────────────────
\paragraph{T5 -- Indépendance entre annales travaillées (V2) et assiduité aux CMs (V14).}

\textbf{Hypothèses.}
$H_0$ : V2 et V14 sont indépendantes ;
$H_1$ : V2 et V14 sont associées.

Les modalités ont été regroupées :
V2 en « 0--1 » ($n=13$), « 2 » ($n=34$), « 3--4+ » ($n=21$) ;
V14 en « Rarement » ($n=6$), « Presque tous » ($n=33$), « Tous » ($n=29$).

\begin{table}[H]
  \centering
  \caption{Tableaux de contingence observé / théorique -- V2 $\times$ V14 ($n = 68$)}
  \label{tab:chi2_T5}
  \small
  \begin{minipage}[t]{0.48\linewidth}
    \centering\textit{Observé}\\[4pt]
    \begin{tabular}{@{}lrrr@{}}
      \toprule
      V2 \textbackslash{} V14 & Rarement & Presque & Tous \\
      \midrule
      0--1   & 0 & 5  & 8  \\
      2      & 3 & 18 & 13 \\
      3--4+  & 3 & 10 & 8  \\
      \bottomrule
    \end{tabular}
  \end{minipage}
  \hfill
  \begin{minipage}[t]{0.48\linewidth}
    \centering\textit{Théorique}\\[4pt]
    \begin{tabular}{@{}lrrr@{}}
      \toprule
      V2 \textbackslash{} V14 & Rarement & Presque & Tous \\
      \midrule
      0--1   & 1{,}15 & 6{,}31 & 5{,}54 \\
      2      & 3{,}00 & 16{,}50 & 14{,}50 \\
      3--4+  & 1{,}85 & 10{,}19 & 8{,}96 \\
      \bottomrule
    \end{tabular}
  \end{minipage}
\end{table}

\textbf{Conditions.}
$n = 68 \geq 50$ {\checkmark} ;
mais $\min(e_{ij}) = 1{,}15 < 5$ {\texttimes}.
La condition sur les effectifs théoriques n'est pas satisfaite pour la
case (0--1, Rarement).
Conformément à la consigne, on présente néanmoins le test à titre indicatif.

\textbf{Critère.}
ddl $= (3-1)(3-1) = 4$ ; seuil $\chi^2_{4\,;\,5\%} = 9{,}49$.
\[
  \chi^2_\text{obs} = \sum_{i,j}\frac{(o_{ij}-e_{ij})^2}{e_{ij}} \approx \mathbf{3{,}61}
\]

\textbf{Conclusion.}
$\chi^2_\text{obs} = 3{,}61 < 9{,}49$ : on \textbf{ne rejette pas $H_0$} (à interpréter
avec précaution au vu de la condition non satisfaite).
Sur cet échantillon, on ne met pas en évidence d'association significative
entre le nombre d'annales travaillées et l'assiduité aux cours magistraux.

% ── T6 ──────────────────────────────────────────────────────────────────────
\paragraph{T6 -- Indépendance entre fréquentation des permanences et moyenne au S1 (SC5/SC6).}

\textbf{Hypothèses.}
$H_0$ : la fréquentation des permanences pédagogiques et la moyenne au S1 sont indépendantes ;
$H_1$ : elles sont associées.

Les données proviennent de \texttt{ADD\_SC5\_SC6.csv} ($n = 79$ paires valides).
La moyenne au S1 a été regroupée en trois classes :
$<10$ ($n=15$), $[10\,;\,12]$ ($n=25$), $>12$ ($n=39$).
La fréquentation des permanences est codée en trois modalités :
« Jamais » ($n=27$), « Rarement » ($n=37$), « Régulièrement » ($n=15$).

\begin{table}[H]
  \centering
  \caption{Tableaux de contingence observé / théorique -- Permanences $\times$ Moy. S1 ($n = 79$)}
  \label{tab:chi2_T6}
  \small
  \begin{minipage}[t]{0.48\linewidth}
    \centering\textit{Observé}\\[4pt]
    \begin{tabular}{@{}lrrr@{}}
      \toprule
      Perm. \textbackslash{} Moy. & $<10$ & $10$--$12$ & $>12$ \\
      \midrule
      Jamais        & 4  & 7  & 16 \\
      Rarement      & 6  & 11 & 20 \\
      Régulièrement & 5  & 7  & 3  \\
      \bottomrule
    \end{tabular}
  \end{minipage}
  \hfill
  \begin{minipage}[t]{0.48\linewidth}
    \centering\textit{Théorique}\\[4pt]
    \begin{tabular}{@{}lrrr@{}}
      \toprule
      Perm. \textbackslash{} Moy. & $<10$ & $10$--$12$ & $>12$ \\
      \midrule
      Jamais        & 5{,}13 & 8{,}54 & 13{,}33 \\
      Rarement      & 7{,}03 & 11{,}71 & 18{,}27 \\
      Régulièrement & 2{,}85 & 4{,}75 & 7{,}41 \\
      \bottomrule
    \end{tabular}
  \end{minipage}
\end{table}

\textbf{Conditions.}
$n = 79 \geq 50$ {\checkmark} ;
$\min(e_{ij}) = 2{,}85 < 5$ {\texttimes}.
La case (Régulièrement, $<10$) ne satisfait pas la condition.
On présente le test à titre indicatif, conformément à la consigne.

\textbf{Critère.}
ddl $= (3-1)(3-1) = 4$ ; seuil $\chi^2_{4\,;\,5\%} = 9{,}49$.
\[
  \chi^2_\text{obs} \approx \mathbf{6{,}73}
\]

\textbf{Conclusion.}
$\chi^2_\text{obs} = 6{,}73 < 9{,}49$ : on \textbf{ne rejette pas $H_0$}
(résultat à interpréter avec prudence).
On ne met pas en évidence d'association significative entre la fréquentation
des permanences pédagogiques et la moyenne au S1 sur cet échantillon.
Ce résultat peut s'expliquer par la faible proportion d'étudiant·es fréquentant
régulièrement les permanences ($n = 15$ sur 79), ce qui limite la puissance du test.
